\section{A generalized reference database for ModelSEED}\label{sec:s6}

\paragraph{A generalized reference database.}
ModelSEED biochemistry is the foundation of ModelSEED, but distinct from its automated pipeline for reconstruction of bacterial metabolism~\cite{henry2010,faria2023}, as well as its counterpart in flowering plant metabolism: PlantSEED~\cite{seaver2014,seaver2018}. Due to the diversity of organisms generated by these pipelines, we strive to maintain an impartiality within this database that would otherwise affect how the models behave.

\paragraph{Transport.} Transport reactions are represented by casting each metabolite into separate compartments, but as many transport reactions may bridge different combinations of compartments, we use agnostic indices: 0 for one compartment and 1 for another. So one reaction can be instantiated against a bacterial membrane, a mitochondrion or a plastid envelope when ModelSEED or PlantSEED reconstructs a model. For this reason, the thermodynamic energy computed for a transport reaction is done so from stoichiometry alone and not from membrane potential or pH gradient.

\paragraph{Reconstruction scope.} The automated reconstruction pipelines in ModelSEED and PlantSEED involve the mapping of biochemistry to enzymatic annotation that is captured in the form of a template, which also includes the custom rules for the assignment of transport. Currently, the union of the ModelSEED v2 and PlantSEED v2 templates houses $\sim9,000$ reactions but this is not a limitation on the scope of the database, and researchers may continue to use any reaction in our database as deemed useful.